% =============================================================================
% Module 1b: Solutions to Exercises
% =============================================================================

\section{Solutions to Exercises}
\label{sec:dg_solutions}

\noindent
Full worked solutions verified in
\solnlink{01b_Differential_Geometry/problems\_01b.wl}{problems\_01b.wl}.

% ----- Exercise 1.1 -----
\subsection*{Exercise~\ref{ex:christoffel_sphere}: Christoffel Symbols of the 2-Sphere}

For $ds^2 = R^2(d\theta^2 + \sin^2\theta\, d\phi^2)$, the metric is
$g_{\theta\theta} = R^2$, $g_{\phi\phi} = R^2\sin^2\theta$, with
$g^{\theta\theta} = 1/R^2$, $g^{\phi\phi} = 1/(R^2\sin^2\theta)$.

Applying~\eqref{eq:christoffel}:
\begin{align*}
    \Gamma^\theta_{\ \phi\phi}
    &= \frac{1}{2} g^{\theta\theta}(-\partial_\theta g_{\phi\phi})
    = \frac{1}{2} \cdot \frac{1}{R^2} \cdot (-2R^2\sin\theta\cos\theta)
    = -\sin\theta\cos\theta \\
    \Gamma^\phi_{\ \theta\phi}
    &= \frac{1}{2} g^{\phi\phi}(\partial_\theta g_{\phi\phi})
    = \frac{1}{2} \cdot \frac{1}{R^2\sin^2\theta} \cdot 2R^2\sin\theta\cos\theta
    = \cot\theta
\end{align*}
All other components vanish (verified by Mathematica).

% ----- Exercise 1.2 -----
\subsection*{Exercise~\ref{ex:christoffel_polar}: Polar Coordinates}

\textbf{(a)} For $ds^2 = dr^2 + r^2 d\phi^2$:
$\Gamma^r_{\ \phi\phi} = -r$, \quad
$\Gamma^\phi_{\ r\phi} = 1/r$.

\textbf{(b)} Computing the Riemann tensor from these Christoffel symbols:
$R^\rho_{\ \sigma\mu\nu} = 0$ for all index combinations.  The space is
flat, confirming that the nonvanishing Christoffel symbols are purely a
coordinate artifact.

\textbf{(c)} The geodesic equations are:
$\ddot{r} - r\dot{\phi}^2 = 0$ and
$\ddot{\phi} + (2/r)\dot{r}\dot{\phi} = 0$.
A straight line $x = a + b\lambda$, $y = c + d\lambda$ in Cartesian
coordinates satisfies these when converted to polar coordinates
$r = \sqrt{x^2 + y^2}$, $\phi = \arctan(y/x)$.

% ----- Exercise 1.3 -----
\subsection*{Exercise~\ref{ex:ricci_sphere}: Curvature of the 2-Sphere}

\textbf{(a)} From~\eqref{eq:riemann}:
$R^\theta_{\ \phi\theta\phi} = \sin^2\theta$.

\textbf{(b)} Ricci tensor: $R_{\theta\theta} = R^\lambda_{\ \theta\lambda\theta}
= R^\phi_{\ \theta\phi\theta} = 1$.
$R_{\phi\phi} = R^\theta_{\ \phi\theta\phi} = \sin^2\theta$.

\textbf{(c)} Ricci scalar:
$R = g^{\theta\theta}R_{\theta\theta} + g^{\phi\phi}R_{\phi\phi}
= \frac{1}{R^2}(1) + \frac{1}{R^2\sin^2\theta}(\sin^2\theta)
= \frac{2}{R^2}$.
This is constant and positive --- the sphere has uniform positive curvature.

% ----- Exercise 1.4 -----
\subsection*{Exercise~\ref{ex:parallel_transport_sphere}: Parallel Transport}

Along a meridian ($\phi = \text{const}$), $d\phi/d\lambda = 0$, so the
parallel transport equation reduces to $dV^\mu/d\lambda = 0$ --- the
vector components are constant along the path.

After traveling north to the pole along $\phi = 0$ and south along
$\phi = \alpha$, the vector returns to the equator.  However, the
coordinate basis $\hat{\phi}$ at $\phi = \alpha$ is rotated by angle
$\alpha$ relative to $\hat{\phi}$ at $\phi = 0$.  The physical vector
has therefore rotated by angle $\alpha$.

Geometrically, the solid angle subtended by the triangular path
(two meridians and a segment of the equator) is $\alpha$.  On a unit
sphere, the holonomy angle equals the enclosed solid angle.

% ----- Exercise 1.5 -----
\subsection*{Exercise~\ref{ex:geodesic_sphere}: Great Circles as Geodesics}

For the equator: $\theta(\lambda) = \pi/2$, $\phi(\lambda) = \lambda$,
so $\dot{\theta} = 0$, $\dot{\phi} = 1$, $\ddot{\theta} = 0$,
$\ddot{\phi} = 0$.

$\theta$-equation: $0 + \Gamma^\theta_{\ \phi\phi} \cdot 1 \cdot 1
= -\sin(\pi/2)\cos(\pi/2) = 0$.  \checkmark

$\phi$-equation: $0 + 2\Gamma^\phi_{\ \theta\phi} \cdot 0 \cdot 1 = 0$.  \checkmark

By the rotational symmetry of the sphere, any great circle can be mapped
to the equator by a rotation, so all great circles are geodesics.
